% Round 18. New-sectors branch. New appendix; independent review required.
\section{Nonlinear-factor monodromy and new full-height families}
\label{r18ng:section}

This appendix treats the geometric points of the deepest
nonlinear strata individually.  Their number changes the
extreme monodromy rank and gives an invariant distinguishing
some of the new generic fibers from every hyperplane model
of HP18.  The local Green input is the separately reviewed
nonlinear-factor theorem NF18.

\subsection{The full geometric top-stratum count}

Use all hypotheses and notation of NF18: $n\ge4$, $m\ge3$,
factor degrees $d_1,\ldots,d_m\ge1$, total degree
$D=\sum_i d_i$, a native smooth normal-crossings arrangement
of the factor hypersurfaces, a degree-$D$ form transverse
to every positive-dimensional stratum and avoiding every
zero-dimensional stratum, and a smooth geometrically
integral generic fiber of $\prod_i f_i+\pi Q$.
Let $\mathcal X/R$ be the explicit NF18 model.
Put
\[
 \mu_I=\prod_{i\in I}d_i,\qquad
 b_n(d_1,\ldots,d_m)
   =\sum_{|I|=n+1}\mu_I-\binom{m-1}{n}.
 \tag{NG18.1}\label{r18ng:rank-number}
\]
An empty sum and a binomial coefficient with upper index
less than its lower index are zero.

\begin{theorem}[Exact top monodromy rank]
\label{r18ng:monodromy}
For every prime $\ell\ne p$, the logarithmic monodromy
operator on $H^n(X_{\overline K},\mathbf Q_\ell)$ satisfies
\[
 \operatorname{rank}N^n=b_n(d_1,\ldots,d_m).
 \tag{NG18.2}\label{r18ng:rank}
\]
More precisely, the extreme terms of the Rapoport--Zink
weight spectral sequence satisfy
\[
 N^n:E_2^{-n,2n}\xrightarrow{\ \sim\ }E_2^{n,0}(-n),
 \qquad
 \dim E_2^{n,0}=b_n(d_1,\ldots,d_m).
 \tag{NG18.3}\label{r18ng:extreme}
\]
This number is positive exactly in the following cases:
$m\ge n+2$, or $m=n+1$ and at least one $d_i>1$.
When it is positive, $X$ has no potential good reduction,
and every strict semistable smooth projective dominant
witness has crossing depth at least $n$.
The rank in \eqref{r18ng:rank} is unchanged by every finite
extension of $K$.
\end{theorem}

\begin{proof}
\emph{Step 1: which strata are connected.}
Every positive-dimensional original complete-intersection
stratum $H_I$ is geometrically connected, and its specified
blowups remain so.  For $|I|=n$ the original stratum is a
smooth geometrically integral curve.  The NF18 blowup
centers on it are empty, because they would be intersections
of $Q$ with zero-dimensional original strata.  Thus this
curve is unchanged.  For $|I|=n+1$, the stratum is the
unchanged smooth zero-dimensional complete intersection
of the factor degrees $d_i$, $i\in I$; by Bezout its
geometric length is exactly $\mu_I$.  Its geometric
points need not be rational over $k$, and none of them
is suppressed in the following complex.

\emph{Step 2: the row of weight zero.}
Let $Y^{[s]}$ be the disjoint union of the $s$-fold
geometric intersections.  The row $b=0$ of the weight
spectral sequence is the alternating restriction complex
$H^0(Y^{[a+1]})$.  In degrees $0\le a\le n-1$, its
basis is the subsets $I$ of cardinality $a+1$, because
all these strata are geometrically connected.  In degree
$n$ its basis has $\mu_I$ points for each subset $I$ of
cardinality $n+1$.  Each such point restricts to the
same list of incident connected curves with the usual
orientation signs.  Thus the last differential is the
ordinary simplex cochain map
\[
 \delta_{n-1}:\mathbf Q^{\binom mn}
                     \longrightarrow\mathbf Q^{\binom m{n+1}}
\]
followed by the injective map that repeats its $I$-coordinate
$\mu_I$ times.  All these statements first make sense as
maps of rational vector spaces on the finite geometric
component sets, with their natural permutation action;
tensoring with $\mathbf Q_\ell$ gives the stated row.

The augmented cochain complex of the full $(m-1)$-simplex
is exact.  Its differential in degree $j$ has rank
$\binom{m-1}{j+1}$: starting with the rank-one augmentation,
this follows inductively from
$\binom m{j+1}-\binom{m-1}j=\binom{m-1}{j+1}$.
The injective coordinate repetition does not change the
rank.  Hence the restriction map into the top geometric
point space has rank $\binom{m-1}n$, and its cokernel
has dimension exactly \eqref{r18ng:rank-number}.
There are no outgoing terms, so this cokernel is
$E_2^{n,0}$.  The convention for $m<n+1$ gives both
spaces zero in this extreme degree.

This argument may also be described using a semisimplicial
dual complex: there is one simplex for every connected
positive-dimensional stratum and $\mu_I$ top simplices
with the same boundary for each $I$ of size $n+1$.
Calling this a single ordinary simplex for each such
$I$ would give an incorrect rank.

\emph{Step 3: the opposite term and the transpose.}
At the other extreme the only $E_1$ summand is
$H^0(Y^{[n+1]},\mathbf Q_\ell)(-n)$; there is no incoming
term.  Its outgoing map is the alternating Gysin map to
$H^2(Y^{[n]},\mathbf Q_\ell)(-(n-1))$.
Each geometric curve in $Y^{[n]}$ is smooth and connected,
so trace identifies its $H^2(1)$ with $\mathbf Q_\ell$.
Pushing forward one of its geometric points has trace
one.  Consequently the Gysin matrix is the transpose,
up to the harmless uniform sign convention, of the
restriction matrix just computed.

Give the rational vector space on all deepest geometric
points the positive definite pairing in which those
points are orthonormal.  It is invariant under the native
Frobenius permutation.  Thus the Gysin kernel is the
orthogonal complement of the restriction image, and
the entire point space is their direct sum.  This is
rational positive-definite linear algebra before tensoring
with $\mathbf Q_\ell$, not an ordering or positivity
claim over $\mathbf Q_\ell$.

For a native closed point of degree $e$, its invariant
fundamental vector is the sum of its $e$ geometric
point vectors; its squared norm and its curve-pushforward
degree are both $e$.  The geometric calculation therefore
retains the full native degrees.  It does not replace
native degree by one for a non-rational closed point.
No analytic descent or assertion that all cycles descend
is used in this cohomology calculation.

\emph{Step 4: the filtered monodromy operator.}
The construction and filtered properties of the weight
spectral sequence are the same primary Saito inputs
verified in HM18.  The map $N^n$ is the identity on the
common extreme $E_1$ point summand, with twist $(-n)$.
It therefore induces the projection of the Gysin kernel
onto the restriction cokernel, which is an isomorphism
by the preceding direct sum.  This proves
\eqref{r18ng:extreme}.  All strata are proper and smooth
over a finite field, so the same purity-of-weights
argument gives degeneration at $E_2$.

For completeness, the weights on the abutment $H^n$
lie in $[0,2n]$.  Indeed a nonzero summand in bidegree
$(a,b)$ with $a+b=n$ has $i\ge\max(0,-a)=\max(0,b-n)$
and $b-2i\ge0$, whence $0\le b\le2n$.
Since $N$ lowers that filtration by two, $N^n$ kills
the piece of weights at most $2n-1$ and has image in
the weight-zero piece.  The latter is its own graded
piece, and \eqref{r18ng:extreme} is an isomorphism
from the top graded piece to it.  Thus its actual rank
is precisely the displayed dimension, proving
\eqref{r18ng:rank} rather than only a lower bound.

Finally, if $m\ge n+1$, then $\mu_I\ge1$ gives
\[
 b_n\ge\binom m{n+1}-\binom{m-1}n
             =\binom{m-1}{n+1}.
\]
This is positive for $m\ge n+2$.  If $m=n+1$, the
number is $\prod_i d_i-1$, positive exactly when
some degree exceeds one.  For $m<n+1$ it is zero.
Finite extension scales normalized logarithmic monodromy
by a nonzero ramification scalar and does not change
the geometric cohomology vector space, hence preserves
this rank.  Potential good reduction and a dominant
witness of depth less than $n$ are excluded by the
already frozen NM17 monodromy argument when the rank
is positive.
\end{proof}

\subsection{A rank distinction from every hyperplane model}

\begin{corollary}[Two explicit fourfold profiles]
\label{r18ng:distinction}
In relative dimension four, five quadratic factors give
total degree ten and
\[
                  \operatorname{rank}N^4=2^5-1=31.
 \tag{NG18.4}\label{r18ng:quadrics}
\]
Five sextic factors give total degree thirty and
\[
                  \operatorname{rank}N^4=6^5-1=7775.
 \tag{NG18.5}\label{r18ng:sextics}
\]
Neither such generic fiber can, even after a finite valued
extension, be isomorphic to a generic fiber of any
four-dimensional hyperplane model HP18 at that same
place, regardless of the hyperplane model's degree.
Both profiles nevertheless have full local ordinary
actual G by NF18.
\end{corollary}

\begin{proof}
The rank formula with $m=n+1=5$ gives the displayed
values.  For a hyperplane model, the exact same rank
formula, with every factor degree one, gives
$\operatorname{rank}N^4=\binom{d-1}5$; it also shows
that HM18's extreme rank is the actual rank.
This sequence is zero for $d\le5$ and strictly increasing
for $d\ge6$.  Its neighboring values around $31$ are
$\binom75=21$ and $\binom85=56$; its neighboring
values around $7775$ are $\binom{17}5=6188$ and
$\binom{18}5=8568$.  Neither required rank occurs.
Rank of logarithmic monodromy is preserved by isomorphism
and finite valued extension, proving the assertion.
The final Green assertion is NF18 in dimension four.
This does not purport to exclude overlap with every
earlier motivic or other proved-G mechanism, or to
exclude all possible dominant witnesses.
\end{proof}

\subsection{The nonlinear arithmetic pencil}

\begin{theorem}[Dense global families for every positive-rank profile]
\label{r18ng:global}
Fix $n\ge4$ and a degree profile $d_1,\ldots,d_m$ for
which $b_n(d_1,\ldots,d_m)>0$.  There is a prime $q$
and a Zariski dense family of rational degree-$D$
hypersurface coefficient classes in $\mathbf P^{n+1}$
with all the following properties.
Every member is smooth and geometrically integral;
at $q$ it has the native NF18 model with those factor
degrees and the exact rank \eqref{r18ng:rank}; at every
other finite place it has the frozen good or
ordinary-quadratic coverage of GP17.  For every number
field $F$ unramified at $q$, its complete Chow groups
carry the height pairing
\[
 \operatorname{CH}^2(X_F)_{\mathbf Q,\mathrm{geom\,hom}}
 \ \times\
 \operatorname{CH}^{n-1}(X_F)_{\mathbf Q,\mathrm{geom\,hom}}
                   \longrightarrow\mathbf R.
 \tag{NG18.6}\label{r18ng:height}
\]
For $n=4$, full local G holds in every codimension and
the complete complementary-degree height pairings
are defined for every $r$.  The prime $q$ may avoid
any prescribed finite set; in particular a family
over any prescribed number field $F$ can be obtained
with $q$ unramified in $F$.
\end{theorem}

\begin{proof}
\emph{Step 1: produce and spread the nonlinear data.}
Over $\mathbf Q$ choose general hypersurface forms
$f_1,\ldots,f_m$ of their prescribed degrees.  The
finite list of all intersection conditions in NF18
is a nonempty open condition: successive Bertini
cuts give smooth complete intersections of every
nonnegative expected dimension, and the intersections
of $n+2$ factors are empty.  Choose a degree-$D$ form
$Q$ transverse to all positive-dimensional members
of this finite list and avoiding the zero-dimensional
ones.  This is another nonempty open Bertini and
avoidance condition.  Rational points exist in these
nonempty opens of coefficient spaces.
After clearing denominators, the data are integral.
The finitely many geometric smoothness, transversality
and avoidance conditions persist modulo every prime
outside a finite set.  Write $P=\prod_i f_i$ and
choose $q>D$, $q\nmid2D$, outside that set and the
prescribed finite set.

\emph{Step 2: impose the exact local congruence.}
Put
\[
 F_0=\sum_{i=0}^{n+1}X_i^D,\qquad
 J_D=\sum_{i=0}^{n}X_i^{D-1}X_{i+1}+X_{n+1}^D,
 \qquad b_D=\prod_{p\mid D}p.
\]
Take $F_1$ in the full affine coefficient lattice
\[
        F_1\equiv P-D F_0+qQ\pmod{q^2},
        \qquad F_1\equiv J_D\pmod{b_D},
 \tag{NG18.7}\label{r18ng:coefficient-congruences}
\]
and in the nonempty fixed-$F_0$ Lefschetz open of
GP17.  This intersection is nonempty and Zariski
dense.  With $F_1$ fixed, choose the finite spreading
set $S$ as in GF18, containing $q$ and every divisor
of $2D$.  Choose the infinite parameter progression
\[
 \begin{split}
 a&\equiv1\pmod{q^2},\qquad a\equiv1\pmod{b_D},\\
 a&\equiv0\pmod p\quad(p\in S,\ p\nmid D,\ p\ne q).
 \end{split}
 \tag{NG18.8}\label{r18ng:parameter-congruences}
\]
Exclude the finitely many singular characteristic-zero
members of $X_a=V(D F_0+aF_1)$.  GP17 supplies smoothness
and geometric integrality of all remaining members.

At $q$ define $Q_a=(D F_0+aF_1-P)/q$.  It is integral
and satisfies $\bar Q_a=\bar Q$, since both congruences
at $q$ were imposed modulo $q^2$.  The exact local
equation is therefore $\prod_i f_i+qQ_a$, with native
uniformizer exponent one and precisely the prescribed
reduced nonlinear data.  NF18 gives codimension-two
G, and Theorem~\ref{r18ng:monodromy} gives its rank.
In dimension four NF18 gives full G.

\emph{Step 3: other places and every new cycle.}
At $p\mid D$ the reduced member is the smooth $J_D$;
at $p\in S$ other than $q$ and not dividing $D$ it
is the smooth Fermat member.  Outside $S$, the
frozen GP17 ordinary-quadratic argument applies with
its allowed arbitrary thickness.  Consequently all
finite places have full codimension-two G.  For a
number-field extension unramified at $q$, the same
native nonlinear model applies anew at each place
above $q$, to every cycle over that completion.
At every other place the good or ordinary-quadratic
theorem is reapplied after base change, with all its
native-degree conventions.  No new cycle is assumed
to be a pullback from the original field.

The complementary one-cycle degree is already
unconditionally covered by the frozen results.
PH16 therefore gives \eqref{r18ng:height} on the
entire displayed groups, and in dimension four gives
all complementary-degree pairings.  GA17 supplies
admissible model lifts on these established domains.

\emph{Step 4: coefficient density with fixed local rank.}
Allow $F_1$ to vary in the fixed lattice and Lefschetz
open in Step 2, making its finite spreading set and
parameter progression only after it is fixed.
If a rational coefficient polynomial $g$ vanishes
on all resulting forms, then for each such $F_1$
the polynomial $g(D F_0+tF_1)$ has infinitely many
integer roots and is identically zero.  At $t=1$
this makes $g$ vanish on the translate of a Zariski
dense lattice-open set.  Hence $g=0$.  This proves
affine coefficient density and therefore projective
coefficient density.  The local factor arrangement,
$q^2$ congruence and monodromy rank are fixed across
this union.  Avoiding the discriminant of a prescribed
number field when choosing $q$ proves the last claim.
\end{proof}

\paragraph{Scope and dependencies.}
For five quadratic or five sextic factors in dimension
four, this theorem gives dense full-height families
whose distinguished local rank is, respectively, $31$
or $7775$.  The rank distinction concerns the actual
generic representation and survives finite extensions;
it is not a change of model of a good generic fiber.
The conclusions do not cover every degree-$D$
hypersurface or arbitrary ramified extensions at $q$,
and no general $U_2$ vanishing is asserted.
The nonlinear geometry and Green theorem are NF18;
the filtered spectral-sequence inputs are the primary
Saito construction cited in HM18; the point count
uses ordinary Bezout and the connectedness input
checked in NF18; the arithmetic and height dependencies
are precisely GP17, PH16 and GA17.  No priority claim
is made for the classical degeneration or pencil methods.
